\section*{Bab \ref{ch:b2:ellingham} --- Diagram Ellingham dan Metalurgi}

\begin{solution}{exo:b2:ellingham:1}
\ce{4Cu + O2 -> 2Cu2O}; \ce{4/3Al + O2 -> 2/3Al2O3}; \ce{2C + O2 -> 2CO};
\ce{C + O2 -> CO2}.
\end{solution}

\begin{solution}{exo:b2:ellingham:2}
$\Delta_r H^\circ = 2(-350.46) = \qty{-700.92}{kJ}$; $\Delta_r S^\circ = 2(43.65)
- 2(41.63) - 205.15 = \qty{-201.1}{J/K}$. Jadi $\Delta_r G^\circ = -700.9 +
0.2011\,T$ (\unit{kJ}, $T$ dalam \unit{K}), berlaku sampai titik lebur seng.
\end{solution}

\begin{solution}{exo:b2:ellingham:3}
Pada \qty{1500}{K}, garis \ce{Cr2O3} berada di dekat \qty{-495}{kJ}. Di
bawahnya, sehingga mampu mereduksi kromium oksida: silikon, titanium,
aluminium, magnesium, kalsium. Di atasnya, tidak mampu: tembaga, besi,
seng. Karbon, pada \qty{-488}{kJ}, terletak sedikit di atasnya: karbon
mereduksi kromium oksida hanya pada suhu yang sedikit lebih tinggi, dan
lalu menghasilkan kromium yang sarat karbida --- salah satu alasan mengapa
kromium untuk paduan bebas karbon dibuat dengan aluminotermi.
\end{solution}

\begin{solution}{exo:b2:ellingham:4}
Kemiringannya adalah $-\Delta_r S^\circ$, yang ditentukan oleh perubahan
jumlah gas. \ce{C + O2 -> CO2}: satu mol gas menghasilkan satu,
$\Delta_r S^\circ \approx 0$, mendatar. \ce{2C + O2 -> 2CO}: satu
menghasilkan dua, $\Delta_r S^\circ > 0$, garisnya turun.
\ce{2CO + O2 -> 2CO2}: tiga menghasilkan dua, $\Delta_r S^\circ < 0$,
garisnya naik, seperti garis-garis logam.
\end{solution}

\begin{solution}{exo:b2:ellingham:5}
$\Delta_r G^\circ(\qty{600}{K}) = -181.6 + 600 \times 0.2165 = \qty{-51.7}{kJ}$,
$p_{\ce{O2}} = p^\circ\exp(-51\,710/(8.314 \times 600)) = \qty{3.2e-5}{bar}$.
Udara mengandung \qty{0.21}{bar}, jauh lebih banyak: raksa teroksidasi pada
\qty{600}{K} (perlahan --- begitulah oksidanya pertama kali dibuat), dan
oksidanya baru terurai di atas sekitar \qty{730}{K}.
\end{solution}

\begin{solution}{exo:b2:ellingham:6}
$\Delta_r G^\circ = -1582.3 - (-743.5) = \qty{-838.8}{kJ}$. Pereaksinya
adalah dua padatan yang hanya bersentuhan pada permukaan butirnya, dan
reaksinya mempunyai energi aktivasi yang tinggi: reaksi itu harus dinyalakan
secara lokal (dengan sumbu magnesium); setelah itu kalor yang
dilepaskannya membuatnya terus berlangsung.
\end{solution}

\begin{solution}{exo:b2:ellingham:7}
Garis \ce{2C + O2 -> 2CO} memotong garis \ce{2Fe + O2 -> 2FeO} di dekat
\qty{1040}{K}; di atasnya, \ce{FeO + C -> Fe + CO} mempunyai \omterm{def:b2:reaction-free-energy:reaction-gibbs}{energi Gibbs
reaksi standar} yang negatif.
\end{solution}

\begin{solution}{exo:b2:ellingham:8}
$\Delta_r G^\circ = 2(-209.1) - (-396.0) = \qty{-22.2}{kJ}$, $K^\circ =
\exp(22\,200/(8.314 \times 1100)) = 11.3$. Maka $x^2/(1 - x) = 11.3$,
$x =
0.92$. Ketika didinginkan sampai \qty{700}{K}, kesetimbangan bergeser
kembali ke arah \ce{CO2} ($x = 0.016$): karbon monoksida terdisproporsionasi,
\ce{2CO -> C + CO2}, dan jelaga mengendap --- perlahan, kecuali jika ada
permukaan seperti besi yang mengkatalisisnya.
\end{solution}

\begin{solution}{exo:b2:ellingham:9}
Garis air (tiga mol gas menghasilkan dua) naik lebih landai daripada garis
\ce{CO}/\ce{CO2}, yang juga tiga menjadi dua tetapi dengan kehilangan
entropi yang lebih besar. Keduanya berpotongan di dekat \qty{1100}{K}: di
bawahnya, garis \ce{CO}/\ce{CO2} lebih rendah dan karbon monoksida adalah
reduktor yang lebih kuat; di atasnya, garis air lebih rendah dan hidrogen
unggul. Ini pernyataan yang sama dengan \omterm{def:b2:equilibrium-shifts:shift}{pergeseran kesetimbangan} gas air
\ce{CO + H2O <=> CO2 + H2} ke kiri pada suhu tinggi.
\end{solution}

\begin{solution}{exo:b2:ellingham:10}
Per mol \ce{O2}, yaitu untuk \ce{2MgO + Si -> SiO2 + 2Mg(g)}: $\Delta_r
G^\circ = -643.7 - (-845.5) = \qty{+201.8}{kJ}$. Dengan silikon dan silika
beraktivitas 1, $\Delta_r G = \Delta_r G^\circ + RT\ln(p_{\ce{Mg}}/p^\circ)^2$,
negatif jika $p_{\ce{Mg}} < p^\circ\exp(-201\,800/(2 \times 8.314 \times 1500)) =
\qty{3e-4}{bar}$. Tanur vakum memompa uap magnesium keluar secepat uap itu
terbentuk dan mengembunkannya pada permukaan dingin, sehingga tekanannya
tidak pernah mencapai nilai itu; kapur yang ditambahkan ke dalam muatan
mengikat silika dan menurunkan aktivitasnya, yang semakin membantu.
\end{solution}

\begin{solution}{exo:b2:ellingham:11}
Empat spesies, satu reaksi, tiga fase (dua padatan dan gas):
$v = 4 - 1 +
2 - 3 = 2$. Pada suhu dan tekanan tertentu, rasio
$p_{\ce{CO}}/p_{\ce{CO2}}$ sudah tertentu: gas itu tidak dapat menyerahkan
seluruh karbon monoksidanya kepada besi oksida, dan gas puncak masih
mengandung banyak karbon monoksida --- yang dibakar untuk memanaskan
udara lebih dahulu.
\end{solution}

\begin{solution}{exo:b2:ellingham:12}
$\log K = 2(0.34 + 0.41)/0.059 = 25.4$, $K = \num{3e25}$. Dengan
$[\ce{Fe^2+}] = \qty{1.0}{mol/L}$, $[\ce{Cu^2+}] = 1/K \approx
\qty{4e-26}{mol/L}$: tembaga tersingkir seluruhnya.
\end{solution}

\begin{solution}{pb:b2:ellingham:1}
\textbf{1.} $\Delta_r H^\circ = \qty{-700.92}{kJ}$, $\Delta_r S^\circ =
\qty{-201.1}{J/K}$.
\textbf{2.} $\Delta_r G^\circ = -700.9 + 0.2011\,T$ (\unit{kJ}), di bawah
\qty{692.7}{K}.
\textbf{3.} $\Delta_{\text{fus}}S = 7320/692.7 = \qty{10.6}{J/(K.mol)}$;
$\Delta_{\text{vap}}S = 115\,300/1180.2 = \qty{97.7}{J/(K.mol)}$.
\textbf{4.} Seng cair: $\Delta_r H^\circ = -700.92 - 2(7.32) = \qty{-715.6}{kJ}$,
$\Delta_r S^\circ = -201.1 - 2(10.6) = \qty{-222.2}{J/K}$: $\Delta_r G^\circ =
-715.6 + 0.2222\,T$.
\textbf{5.} Seng gas: $\Delta_r H^\circ = -715.6 - 2(115.3) = \qty{-946.2}{kJ}$,
$\Delta_r S^\circ = -222.2 - 2(97.7) = \qty{-417.7}{J/K}$: $\Delta_r G^\circ =
-946.2 + 0.4177\,T$.
\textbf{6.} Kemiringan $0.201$, $0.222$, dan $0.418$ \unit{kJ/K}. Di atas
titik didih, reaksi mengonsumsi tiga mol gas, bukan satu: kehilangan
entropinya, dan dengan demikian kemiringannya, menjadi dua kali lipat.
\textbf{7.} Pada \qty{692.7}{K}: $-700.9 + 139.3 = -715.6 + 153.9 = \qty{-561.6}{kJ}$;
pada \qty{1180.2}{K}: $-715.6 + 262.3 = -946.2 + 492.9 = \qty{-453.3}{kJ}$.
Garis-garisnya tersambung, sebagaimana mestinya.
\textbf{8.} Kemiringan $(-453.0 + 418.2)/200 = \qty{-0.174}{kJ/K}$, $a = -418.2 +
0.174 \times 1100 = \qty{-226.8}{kJ}$: $\Delta_r G^\circ = -226.8 - 0.174\,T$.
\textbf{9.} $\Delta_r S^\circ = \qty{+174}{J/K}$: satu mol gas menghasilkan
dua.
\textbf{10.} \ce{2ZnO + 2C -> 2Zn(g) + 2CO}, $\Delta_r G^\circ = (-226.8 - 0.174\,T)
- (-946.2 + 0.4177\,T) = 719.4 - 0.592\,T$ (\unit{kJ}).
\textbf{11.} Nol pada $T = 719.4/0.592 = \qty{1215}{K}$.
\textbf{12.} Di atas titik didihnya, seng terbentuk sebagai uap: reaksi itu
menghasilkan gas dari padatan, entropinya besar dan positif, dan uapnya
meninggalkan muatan, yang memisahkan logam dari bijih dan kokas dengan
distilasi.
\textbf{13.} Empat spesies, satu reaksi, satu syarat khusus
($p_{\ce{Zn}} =
p_{\ce{CO}}$), tiga fase: $v = 4 - 1 - 1 + 2 - 3 = 1$.
Menetapkan tekanan berarti menetapkan suhu kesetimbangan.
\textbf{14.} Satu mol uap seng per mol karbon monoksida:
$p_{\ce{Zn}} = p_{\ce{CO}} = \qty{0.5}{bar}$.
\textbf{15.} Per mol \ce{ZnO}, $\Delta_r G^\circ = \frac12(719.4 - 0.592 \times
1300) = \qty{-25.1}{kJ/mol}$, $K^\circ = \exp(25\,100/(8.314 \times 1300)) = 10$;
$Q = 0.25 < K^\circ$: reduksi berjalan.
\textbf{16.} Tepat di atas perpotongan, reduksi sudah diuntungkan (dengan
\qty{0.5}{bar} masing-masing gas, mulai sekitar \qty{1170}{K}); suhu yang
lebih panas memboroskan bahan bakar, merusak retort tanah liat, dan tidak
menghasilkan seng lebih banyak, karena reaksinya toh berlangsung tuntas.
\textbf{17.} \ce{Zn(g) + CO2 -> ZnO + CO}.
\textbf{18.} Garis seng oksida terletak di bawah garis \ce{CO}/\ce{CO2}
($-445$ dibandingkan \qty{-357}{kJ} pada \qty{1200}{K}): uap seng
mereduksi karbon dioksida dan kembali menjadi oksida. Di dalam retort,
karbon panas memusnahkan setiap \ce{CO2} (Boudouard); di dalam kondensor
yang lebih dingin tidak ada yang melakukannya, dan \ce{CO2} --- dari udara
yang bocor masuk, atau dari \ce{2CO -> C + CO2} ketika pendinginan ---
merusak seng menjadi serbuk kelabu yang teroksidasi.
\textbf{19.} Pendinginan lambat membiarkan uap itu lama berada dalam
rentang tempat uap teroksidasi kembali, dan menghasilkan debu halus alih-alih
logam cair; pengembunan cepat menjadi cairan membatasi keduanya.
\textbf{20.} $10^6/65.38 = \qty{1.53e4}{mol}$ seng, dan karbon sebanyak
itu pula: $1.53 \times 10^4 \times 12.011 = \qty{184}{kg}$ (dalam praktik
jauh lebih banyak, untuk memanaskan retort).
\textbf{21.} $\Delta_r H^\circ = \frac12(-226.8 + 946.2) = \qty{+360}{kJ}$
per mol seng; per ton, $1.53 \times 10^4 \times 360 = \qty{5.5e6}{kJ}$,
yaitu \qty{5.5}{GJ} --- disediakan dengan membakar lebih banyak batu bara
di sekeliling retort.
\textbf{22.} Karbon mereduksi seng oksida pada kondisi standar di atas
$\boldsymbol{T \approx \qty{1.22e3}{K}}$.
\end{solution}
